\section{Real data: what exact attention shows in a language model}\label{sec:real}

\subsection{Data and validation}
The real tensors come from Qwen2.5-0.5B~\cite{qwen}, an open language model with 24
layers, $H=14$ query heads that share $H_{kv}=2$ key/value heads ($G=7$), and head
dimension $d=d_v=64$, which is also the head dimension of the synthetic runs (the
grouped-head shape is used in \cref{sec:gqa}; the scaling runs have one head). The input
is the first $N$ tokens of a long public-domain book (\emph{War and Peace}, Project
Gutenberg~\#2600, 773{,}025 tokens in total) for $N=8{,}192$, $16{,}384$ and
$32{,}768$, the largest context the model supports. A Python script runs the model
once on a GPU and records, for layers 0, 11 and 23, the queries and keys \emph{after}
the model's rotary position embedding, the values, and the attention output before
the output projection. Python is used for this extraction and, elsewhere, for
submitting jobs and drawing figures; all attention results below are recomputed by
the C++ and CUDA code from these files.

\inputifexists{report/generated/table_v7.tex}

\Cref{tab:v7} compares three computations of the same quantity. My distributed
single-precision result agrees with my double-precision reference to a few
$10^{-5}$ of the output scale, and with PyTorch's own single-precision output to
the same level; the GPU kernel behaves alike. The measure is the largest absolute
difference divided by the largest output component of the tensor; it is not a
relative error per element, and a head whose output is small could hide behind a
large one. The outputs of the real layers are not of order one (their largest
component ranges from \n{real.scale.min} to \n{real.scale.max} across the nine
tensors), so the automatic check of a run replaces the fixed tolerance $10^{-5}$ of
the synthetic outputs by $10^{-3}$ times the largest component of the reference rows
it compares with (the \n{real.valrows} sampled rows; their largest component is at
most that of the whole tensor). That bound was written into the job script before
the first real-tensor run and was not changed afterwards. It is deliberately
loose, a guard against gross errors in the head mapping or the mask and not a
statement about accuracy: every measured difference is at least
\n{real.tol.margin} times smaller, and the measured values in \cref{tab:v7} are the
result to read.
All seven communication variants were also run on one layer with every row
checked; they agree with the reference to the same level.

One lesson from this step concerns validation itself. The first version of the
extraction script checked its own dump with a short formula and reported errors of
order one on some rows. The dump was correct: PyTorch's attention function,
evaluated in double precision on the dumped tensors, reproduced the model's
single-precision output on every row. The disagreement was in my check, where an
\code{einsum} contraction and explicit sums gave different values on those rows. I
did not trace the cause (precision, memory layout or backend), so I make no claim
about the library; the script now uses explicit sums and reports the comparison
with PyTorch's double-precision attention. An independent check is only evidence
if it is itself checked.

\subsection{Grouped-query attention moves seven times fewer bytes}\label{sec:gqa}
With GQA only the two key/value heads travel, while the work is that of fourteen
query heads. The matched control replicates each key/value head seven times, which
leaves the mathematics unchanged (the outputs are bit-identical in the unit test)
and multiplies the traffic by exactly seven.

With 72 ranks on one node and synthetic tensors of the model's shape, the
overlapped ring needs \n{gqa.T.overlap.N32768.P72}\,s with GQA and
\n{gqa.Tmha.overlap.N32768.P72}\,s with the control at $N=32{,}768$
(\n{gqa.gain.pct.overlap.N32768.P72}\% more), and a rank receives
\n{gqa.mib.N32768.P72}\,MiB instead of \n{gqa.mib.mha.N32768.P72}\,MiB, a ratio of
\n{gqa.bytes.ratio.N32768.P72}. The network is not what makes the difference here,
however: the control M0, which exchanges nothing, is slower by almost the same
amount (\n{gqa.gain.pct.local.N32768.P72}\%). The seven query heads of a group read
the same packed key tile and the same values; with GQA the kernel therefore streams
one seventh of the key/value data through the caches, which is the likely cause. The real tensors behave alike
($N=32{,}768$, layer 11, causal mask, zigzag partition, 72 ranks):
\n{gqa.real.T.overlap.P72}\,s against \n{gqa.real.Tmha.overlap.P72}\,s for the ring
and \n{gqa.real.T.local.P72}\,s against \n{gqa.real.Tmha.local.P72}\,s for the
control. The third prediction of \cref{sec:model} was therefore only half right:
GQA does divide the bytes by $G$, and it also speeds up the local kernel, which
the model does not contain.

The network term of the prediction appears where the model says it should, when
ranks own few rows and the exchange crosses nodes. With 288 ranks on four nodes and
$N=8{,}192$ (28 rows per rank) the control is \n{gqa.gain.pct.overlap.N8192.P288}\%
slower than GQA with the overlapped ring and \n{gqa.gain.pct.ring.N8192.P288}\% slower
with the blocking ring, while M0 differs by \n{gqa.gain.pct.local.N8192.P288}\% as
before; at $N=32{,}768$ (114 rows per rank) the three figures are
\n{gqa.gain.pct.overlap.N32768.P288}\%, \n{gqa.gain.pct.ring.N32768.P288}\% and
\n{gqa.gain.pct.local.N32768.P288}\%.

\subsection{How far back does the model look?}
Exact attention over the full context is the ground truth against which cheaper
schemes can be judged. For 2{,}048 evenly spaced query positions of the
$N=32{,}768$ context and every head, a separate double-precision two-pass program
(\code{attn\_analyze}) computes the exact probabilities and derives: the probability mass that falls on the most
recent $w$ tokens; the relative error $\|O_w-O\|/\|O\|$ of the output when attention
is restricted to that window and renormalised (the \emph{window} approximation);
and the same error when the first four tokens of the sequence are kept as well (the
\emph{sink+window} approximation used by streaming language models~\cite{streaming}).
The statistics below use the queries in the second half of the sequence, whose
context is longer than every window considered. This analysis needs no exchange
between ranks: every rank reads all keys and values of the layer and takes its
share of the sampled queries. The distributed ring is not what produces these
statistics; its role on the real tensors is the validation of \cref{tab:v7} and the
timings of \cref{sec:gqa}. What is compared is the output of
one attention layer for fixed inputs. This is a replay of single layers, not a run
of the model with windowed attention in every layer (where the change in one layer
would alter the inputs of the next), so the errors below are not losses of model
quality and say nothing about perplexity or task accuracy.

\begin{figure}[t]
\centering
\figureifexists{real_window.pdf}
\caption{Qwen2.5-0.5B, $N=32{,}768$, queries at positions $\ge 16{,}384$, all 14
heads. Lines are medians over (query, head) pairs and bands the interquartile range
of those pairs (a description of their distribution, not a confidence interval:
neighbouring tokens and heads are not independent samples).
(a) Probability mass on the most recent $w$ tokens. (b) Relative output error of the
window approximation. (c) The same with the first four tokens kept.}
\label{fig:realwindow}
\end{figure}

\begin{figure}[t]
\centering
\figureifexists{real_heads.pdf}
\caption{The same data per head (medians). (a) Error of a 1{,}024-token window; the
colour saturates at 1. (b) Probability on the first token. Heads 0--6 read
key/value head 0 and heads 7--13 read key/value head 1.}
\label{fig:realheads}
\end{figure}

\paragraph{Three layers, three behaviours.}
\Cref{fig:realwindow,fig:realheads} show that the answer depends strongly on depth.
\begin{itemize}
\item \textbf{Layer 0 is local.} The median query places \n{real.mass_w.layer0.1024}\% of
its probability on the last 1{,}024 tokens and \n{real.mass_w.layer0.4096}\% on the last
4{,}096. A 1{,}024-token window changes the median output by
\n{real.err_w.layer0.1024}\% and a 4{,}096-token window by
\n{real.err_w.layer0.4096}\%. Heads differ, however: \n{real.heads.exact.layer0} of the
fourteen are reproduced essentially exactly by a 1{,}024-token window and the others
lose \n{real.heads.rest.min.layer0}--\n{real.heads.rest.max.layer0}\%.
\item \textbf{Layer 11 has attention sinks.} \n{real.sinkheads.count} of the fourteen
heads, precisely those that share key/value head 1 (heads
\n{real.sinkheads.first}--\n{real.sinkheads.last}), put between \n{real.sinkheads.min}\%
and \n{real.sinkheads.max}\% of their probability on the very first token of the
sequence, 16{,}000 to 32{,}000 positions away; no other head exceeds
\n{real.otherheads.max}\%. A window shorter than the distance to the start of the sequence removes that
token; the remaining probabilities are then renormalised and the output changes
completely (the median
relative error is \n{real.err_w.layer11.1024} for $w=1{,}024$, and one (query, head)
pair in ten has an error above \n{real.err_w.p90.layer11.1024}, because the exact
output of a sink head is small and the renormalised window output is not). Keeping
four sink tokens repairs this: the median error falls to
\n{real.err_sw.layer11.1024}\% for $w=1{,}024$, \n{real.err_sw.layer11.4096}\% for
$w=4{,}096$ and \n{real.err_sw.layer11.16384}\% for $w=16{,}384$.
\item \textbf{Layer 23 is diffuse.} Its median entropy is \n{real.entropy.layer23}
nats, the entropy of a uniform distribution over about \n{real.entropy.eff.layer23}
tokens, and there is no sink. \n{real.beyond.layer23.1024}\% of the probability lies
beyond the last 1{,}024 tokens and \n{real.beyond.layer23.4096}\% beyond the last
4{,}096. A 1{,}024-token window changes the median output by
\n{real.err_w.layer23.1024}\%, a 4{,}096-token window by \n{real.err_w.layer23.4096}\%,
and even a 16{,}384-token window by \n{real.err_w.layer23.16384}\%.
\end{itemize}

\paragraph{Conclusions about the data.}
For this model, this text and these three layers: (i) for the queries of the second
half and windows of up to 16{,}384 tokens, a pure sliding window fails in the middle
layer, whose sink heads depend on the start of the sequence; (ii) keeping four
sink tokens removes that failure but still changes the output of the last layer by
a third at a window of 1{,}024 tokens; and (iii) the sink heads are exactly the
heads of one key/value group. The last point is an observation; it does not show
that grouping causes sinks, and a high attention weight is not by itself evidence
that a token matters for the result. These statements are descriptive. They rest
on one model, one text and three layers, and the (query, head) pairs behind each
median are correlated, so they show what the exact computation can measure, not a
general law. For the purpose of this project they answer the motivating question:
at a 32{,}768-token context a substantial share of the attention of a real model
is long-range, so an exact method has something to compute that windows discard,
and the distributed scheme studied here is how the exact output of such a layer
is obtained for all queries when one process is not enough.
